\justifying
\NSFCBodyText

\subsubsubsection{研究内容}

本项目面向新疆口岸—中亚运输—边境仓储链路中观测不完备、网络间歇与端侧资源受限并存的条件，围绕“货物可靠感知—弱网资源协同—仓储风险预警”开展三项递进且可独立验证的研究。

\textbf{（1）货物多模态智能感知与异常识别。}融合振动、倾斜、电子封志、光感和包装图像，研究不同采样频率下的时间对齐、观测质量估计和跨模态可靠性自适应表征。建立运输全过程的事件状态模型，识别异常冲击/倾斜、非法开启以及包装破损、变形等异常，并输出事件概率、损伤位置和预测不确定性。通过模态删除、时间错位、信号噪声、图像遮挡及跨任务留出实验，检验互补证据在何种条件下仍能支持可信判断。

\textbf{（2）弱网通信、边缘计算与低功耗协同。}面向 2G/3G、高丢包、低带宽、长时延和间歇连接，将网络状态、缓存队列、可用算力、剩余电量与事件风险纳入统一决策状态。依据风险/信息价值联合选择本地推理深度、自适应压缩、离线缓存、断点续传、优先上传和动态工作周期，使高风险证据在弱网中可靠送达，低价值数据在端侧摘要或延迟传输；在相同事件流和网络轨迹上揭示可靠性、时延、通信量与能耗之间的协同边界。

\textbf{（3）仓储环境多源感知与异常预警。}融合温湿度、安防视频、烟雾和门禁信息，研究缓慢环境变化、突发安防事件及其跨源共现关系，建立可靠性加权的仓储风险状态演化模型。形成由实时监测、复合异常识别、校准风险评估到分级预警和应急响应建议的闭环。运输阶段货物状态在具备永久货物标识和前瞻配对记录时作为可选先验；缺少跨阶段数据不影响仓储四类信息的独立研究与验证。

\subsubsubsection{研究目标}

总体目标是建立适用于弱网跨境物流的多模态风险感知与边缘协同方法体系，使运输货物状态、端侧资源动作和仓储环境风险在统一的事件级评价与不确定性口径下可建模、可验证、可证伪。具体目标为：一是形成面向五类运输观测的可靠性自适应融合方法，实现全过程异常识别、包装损伤定位和可信告警；二是形成覆盖边缘推理、压缩、缓存、断点续传、可靠上传和低功耗管理的联合策略，在安全约束下验证资源收益；三是形成面向温湿度、视频、烟雾和门禁的仓储复合风险演化与分级预警方法，验证其提前量、误报控制、校准性和跨条件稳健性。

\subsubsubsection{拟解决的关键科学问题}

\textbf{（1）不完备观测下货物多模态互补信息的可靠表征与异常辨识。}在异步、缺失、噪声和可靠性漂移共同存在时，哪些跨模态互补关系仍能稳定表征货物状态；如何区分真实冲击、非法开启、包装破损/变形与观测质量下降，并避免缺失模态造成置信度虚高？

\textbf{（2）通信--计算--能量耦合约束下的信息价值分配与边缘协同边界。}风险程度、信息新颖度、按时送达概率、缓存积压、端侧算力和剩余电量如何共同决定感知、推理、压缩、缓存、上传/续传及工作周期，使高风险证据可靠送达并形成可解释的可靠性--时延--通信量--能耗边界？

\textbf{（3）异质仓储多源证据驱动的复合风险演化与提前预警。}不同时标、不同可靠性的温湿度、视频、烟雾和门禁证据如何形成统一风险状态；哪些跨源共现、持续性和先后关系能够区分环境异常、安防异常与复合风险，并在控制误报的同时获得可校准的分级预警？
